\section{Mathematical Formulation}
\label{sec:formulation}

% =============================================================================
% 2.1  Geometry and coordinate system
% =============================================================================
\subsection{Oblate spheroidal shell model}
\label{sec:geometry}

We model the watermelon as an oblate spheroid with semi-major axis $a$
(equatorial radius), semi-minor axis $c$ (polar radius), and uniform
rind thickness $h$.  The aspect ratio is defined as
$\zeta = c/a$,
where $\zeta < 1$ corresponds to the characteristic oblate shape of most
commercial watermelon cultivars ($\zeta \approx 0.5$--$0.7$ for
elongated varieties; $\zeta \approx 0.8$--$0.9$ for round cultivars).
The equivalent-volume sphere has radius
\begin{equation}
  R_{\mathrm{eq}} = \left( a^2 c \right)^{1/3}.
  \label{eq:Req}
\end{equation}

The rind is treated as a thin, isotropic, viscoelastic shell with Young's
modulus $E_{\mathrm{rind}}$, Poisson's ratio $\nu$, density
$\rho_w$, and structural loss factor $\eta$.  The enclosed flesh is
modelled as an inviscid, nearly incompressible fluid with density
$\rho_f$ and bulk modulus $K_f$.  The canonical parameter values used
throughout this study are given in Table~\ref{tab:canonical}.

\begin{table}[t]
  \centering
  \caption{Canonical parameters for the Crimson Sweet watermelon at
    optimal ripeness.  Values are drawn from Sadrnia et
    al.~\cite{Sadrnia2006}, De Belie et al.~\cite{DeBelieetal2000},
    and established tissue mechanics compilations.}
  \label{tab:canonical}
  \begin{tabular}{llcc}
    \toprule
    Parameter & Symbol & Value & Unit \\
    \midrule
    Semi-major axis    & $a$                 & 0.158 & m \\
    Semi-minor axis    & $c$                 & 0.123 & m \\
    Aspect ratio       & $\zeta = c/a$       & 0.78  & --- \\
    Equivalent radius  & $R_{\mathrm{eq}}$   & 0.145 & m \\
    Rind thickness     & $h$                 & 15    & mm \\
    Rind modulus (ripe) & $E_{\mathrm{rind}}$ & 50   & MPa \\
    Poisson's ratio    & $\nu$               & 0.40  & --- \\
    Rind density       & $\rho_w$            & 1050  & kg/m$^3$ \\
    Flesh density      & $\rho_f$            & 950   & kg/m$^3$ \\
    Flesh bulk modulus  & $K_f$              & 2.2   & GPa \\
    Turgor pressure     & $P_{\mathrm{int}}$ & 500   & Pa \\
    Loss tangent        & $\eta$             & 0.15  & --- \\
    \bottomrule
  \end{tabular}
\end{table}

% Schematic figure (placeholder — no diagram generated yet)
% \begin{figure}[t]
%   \centering
%   \includegraphics[width=0.6\textwidth]{fig_watermelon_schematic}
%   \caption{Geometry of the oblate spheroidal shell model.  The
%     watermelon is characterised by semi-major axis~$a$, semi-minor
%     axis~$c$, rind thickness~$h$, and flesh properties
%     ($\rho_f$, $K_f$).  The coordinate system and modal
%     decomposition are indicated.}
%   \label{fig:schematic}
% \end{figure}

% =============================================================================
% 2.2  Eigenfrequency derivation
% =============================================================================
\subsection{Eigenfrequency of the flexural mode}
\label{sec:eigenfreq}

For a thin spheroidal shell filled with fluid, the natural frequency of
the $n$th flexural mode may be written in the general form
\begin{equation}
  f_n = \frac{1}{2\pi}
  \sqrt{
    \frac{E_{\mathrm{rind}}\, h}{\rho_{\mathrm{eff}}\, R_{\mathrm{eq}}^3}
    \, \mathcal{G}_n(\nu, \zeta)
  },
  \label{eq:fn_general}
\end{equation}
where $\mathcal{G}_n(\nu, \zeta)$ is a dimensionless geometric factor
that depends on the mode number~$n$, Poisson's ratio~$\nu$, and the
aspect ratio~$\zeta$, and $\rho_{\mathrm{eff}}$ is an effective
density incorporating the fluid loading,
\begin{equation}
  \rho_{\mathrm{eff}} = \rho_w h + \rho_f\, R_{\mathrm{eq}}\,
    \mathcal{F}_n(\zeta),
  \label{eq:rho_eff}
\end{equation}
where $\mathcal{F}_n(\zeta)$ is the fluid-loading shape factor for
mode~$n$ \cite{Junger1986, lamb1882}.

The dominant tap-tone mode is the lowest flexural mode, $n = 2$.  For
this mode, the squared frequency is
\begin{equation}
  f_2^2 = \frac{E_{\mathrm{rind}}\, h\, \mathcal{G}_2(\nu, \zeta)}
    {4\pi^2\, \rho_{\mathrm{eff}}\, R_{\mathrm{eq}}^3}.
  \label{eq:f2_squared}
\end{equation}

\noindent A key observation: $f_2^2$ is \emph{exactly linear} in
$E_{\mathrm{rind}}$ when turgor pressure is negligible.  This linearity
is the foundation of the closed-form inversion developed in
\S\ref{sec:inversion}.

The geometric factor~$\mathcal{G}_2$ incorporates both bending and
membrane stiffness contributions.  In the Rayleigh--Ritz framework
\cite{Leissa1973, Soedel2004}, the modal stiffness for mode~$n$ of a
thin shell decomposes as
\begin{equation}
  K_n = K_{\mathrm{bend}} + K_{\mathrm{memb}} + K_{\mathrm{prestress}},
  \label{eq:stiffness_decomp}
\end{equation}
where the bending stiffness scales as
\begin{equation}
  K_{\mathrm{bend}} = \frac{n(n-1)(n+2)^2\, h^3}
    {12(1-\nu^2)\, R_{\mathrm{eq}}^4}\, E_{\mathrm{rind}},
  \label{eq:K_bend}
\end{equation}
the membrane stiffness scales as
\begin{equation}
  K_{\mathrm{memb}} = \frac{h}{R_{\mathrm{eq}}^2}\,
    \frac{n^2 + n - 2 + 2\nu}{n^2 + n + 1 - \nu}\,
    E_{\mathrm{rind}},
  \label{eq:K_memb}
\end{equation}
and the pre-stress stiffness (from internal turgor pressure~$P_{\mathrm{int}}$) is
\begin{equation}
  K_{\mathrm{prestress}} = \frac{P_{\mathrm{int}}}{R_{\mathrm{eq}}}
    \, (n-1)(n+2).
  \label{eq:K_prestress}
\end{equation}

For the canonical ripe watermelon ($n = 2$, $h = \SI{15}{\milli\metre}$,
$R_{\mathrm{eq}} = \SI{0.145}{\metre}$), the bending and membrane
terms both scale linearly with~$E_{\mathrm{rind}}$, whereas the
pre-stress term is independent of~$E$.  This separation underpins the
linearity $f_2^2 = \alpha E_{\mathrm{rind}}$, where the proportionality
constant is
\begin{equation}
  \alpha = \frac{h\,\mathcal{G}_2(\nu, \zeta)}
    {4\pi^2\, \rho_{\mathrm{eff}}\, R_{\mathrm{eq}}^3}
    = \SI{161.5}{\hertz\squared\per\mega\pascal}
  \label{eq:alpha_value}
\end{equation}
for the canonical geometry.

% =============================================================================
% 2.3  Closed-form inversion
% =============================================================================
\subsection{Inversion: from tap-tone to rind modulus}
\label{sec:inversion}

Rearranging Eq.~\eqref{eq:f2_squared} for $E_{\mathrm{rind}}$ gives
\begin{equation}
  \boxed{
    E_{\mathrm{rind}} = \frac{4\pi^2\, \rho_{\mathrm{eff}}\,
      R_{\mathrm{eq}}^3}{h\, \mathcal{G}_2(\nu, \zeta)}\, f_2^2.
  }
  \label{eq:inversion}
\end{equation}
Thus a single measurement of the tap-tone frequency~$f_2$, combined
with knowledge (or estimation) of the fruit geometry and flesh density,
yields the rind elastic modulus directly.  No iterative solver or
optimisation is required.

More explicitly, the inversion takes the form
\begin{equation}
  E_{\mathrm{rind}} = \frac{\omega^2\, m_{\mathrm{eff}}
    - K_{\mathrm{prestress}}}
    {K_{\mathrm{bend}}^{(0)} + K_{\mathrm{memb}}^{(0)}},
  \label{eq:inversion_explicit}
\end{equation}
where $\omega = 2\pi f_2$, $m_{\mathrm{eff}} = \rho_w h + \rho_f
R_{\mathrm{eq}} / n$ is the effective mass per unit area including
fluid loading, and $K_{\mathrm{bend}}^{(0)}$ and
$K_{\mathrm{memb}}^{(0)}$ are the geometric pre-factors of
$E_{\mathrm{rind}}$ in Eqs.~\eqref{eq:K_bend}--\eqref{eq:K_memb}.
Equation~\eqref{eq:inversion_explicit} is algebraically exact: the
round-trip error $|E_{\mathrm{inv}} - E_{\mathrm{true}}| /
E_{\mathrm{true}}$ is less than \SI{0.01}{\percent} for all ripeness
stages tested (\S\ref{sec:forward}).

% =============================================================================
% 2.4  Dimensional analysis and the ripeness parameter
% =============================================================================
\subsection{Dimensionless ripeness parameter}
\label{sec:Pi_ripe}

To collapse predictions across cultivars, we define the dimensionless
ripeness parameter
\begin{equation}
  \Pi_{\mathrm{ripe}} = f_2 \, R_{\mathrm{eq}}
    \sqrt{\frac{\rho_{\mathrm{eff}}}{E_{\mathrm{rind}}}},
  \label{eq:Pi_ripe}
\end{equation}
where $\rho_{\mathrm{eff}} = \rho_w h / R_{\mathrm{eq}} + \rho_f$.
From Eq.~\eqref{eq:f2_squared}, this parameter evaluates to a universal
constant
\begin{equation}
  \Pi_{\mathrm{ripe}} = \frac{1}{2\pi}
    \sqrt{h\, \mathcal{G}_2(\nu, \zeta) / R_{\mathrm{eq}}^3}
  \label{eq:Pi_ripe_universal}
\end{equation}
for geometrically similar fruits.  Deviations from the master curve
indicate cultivar-specific geometry or constitutive effects.

For the four cultivars considered in this study (Crimson Sweet, Sugar
Baby, Charleston Gray, and Seedless Tri-X~313),
$\Pi_{\mathrm{ripe}} = 0.062 \pm 0.002$ (CoV~$= \SI{3.3}{\percent}$;
see Table~\ref{tab:Pi_ripe} in \S\ref{sec:universal_curve}).  This
tight collapse confirms that the dimensionless parameter effectively
normalises out the dominant cultivar-specific geometric variation.
